\documentclass[a4paper,10pt]{report}

%picture imports
\usepackage{graphicx}
\usepackage{color}
\definecolor{lbcolor}{rgb}{0.85,0.85,0.85}

%sourcecode formatting
\usepackage[final]{listings}
\usepackage{url}


\lstloadlanguages{Eiffel,bash}
\newcommand{\eiffellisting}{\lstset{language=Eiffel,tabsize=2,backgroundcolor=\color{lbcolor},frame=tb,breaklines=true}}
\newcommand{\bashlisting}{\lstset{basicstyle=\small\tt,language=bash,tabsize=2,backgroundcolor=\color{lbcolor},frame=tb,breaklines=true,morekeywords={svn,geant,mkdir,activate_estudio}}}
\newcommand{\identifier}[1]{\texttt{#1}}

\eiffellisting
\begin{document}

%opening
\title{Selective Capture and Replay for Eiffel}
\author{Stefan Sieber}

\begin{titlepage}
	\begin{flushleft}
		\includegraphics[width=100pt]{illustrations/SE.png}
	\end{flushleft}

	\vskip 50pt
	\begin{center}
		\Huge \textbf{Capture and Replay for Eiffel}
	\end{center}

	\vskip 100pt
	\begin{flushright}
		\large \textbf{Stefan Sieber}\\
		\vskip 10pt
		02-911-790\\
		Master Thesis\\
		Summer 2007\\
	\end{flushright}
	
	\vskip 20pt
	\begin{flushright}
		\large Chair of Software Engineering\\
		Department of Computer Science\\
		ETH Z\"{u}rich
	\end{flushright}
		
	\vskip 20pt
	\begin{flushright}
		\large Andreas Leitner\\
		Prof. Bertrand Meyer
	\end{flushright}

	\vskip 50pt
	\begin{tabular*}{1.00\textwidth}{@{\extracolsep{\fill}}lr}
		 \includegraphics[width=100pt]{illustrations/ETH.png} &
		 \includegraphics[width=100pt]{illustrations/INF.png}
	\end{tabular*}
	\vfil
\end{titlepage}

\newpage
\begin{abstract}
%Should the motivation include a reference to cdd?
Debugging applications is very tedious work. One of the hardest steps to find a fault in a program is to reproduce the steps that lead to the associated program failure. Capture and replay aims at simplifying this process. When trying to replay a software it is necessary to decide what parts are to be simulated in order to drive the rest of the program. 

Selective Capture and Replay \cite{orso05may} makes it possible to replay the same program run again on a predefined set of classes, thus reproducing those steps automatically. It is possible to define the subset of the program that should be replayed individually. This makes it special compared to conventional systems where this border is fixed (e.g. everything except the network library). Another unique feature is that the amount of recorded data is linear to the number of variables passed over this border, whereas conventional systems record all the data that passes the border.
 The goal of this thesis is to implement Selective Capture Replay for Eiffel, which in contrast to Java, supports multiple inheritance. This makes it necessary to instrument the programs in a different way than in the original implementation. One benefit of this changed implementation is the possibility to switch between capture and replay phase without recompilation.
 To validate the implementation, an example will be presented and the performance of the technique will be measured using this example.

%CDD [reference!] facilitates this step by providing the developer a testcase that results in the same failure. 
%motivation
%capture/replay
%scope
%-modifications
%-implementation
%-experimental results
\end{abstract}

\newpage

\section{Introduction}

\chapter{Related Work}
%TODO: kurze einleitung

\section{Overview}
Because replaying application runs is important in order to be able to debug or test applications, there exist many techniques that implement capture and replay. One of the basic steps in order to be able to capture and replay an application is to distinct between the applications deterministic core (the \emph{observed part}) and the non-deterministic environment like user input, network or external storage (the \emph{unobserved part}). Capture and replay techniques capture the interactions between observed and unobserved part during \emph{capture phase} so that they can replay these interactions on the observed part during the \emph{replay phase}.

% -------
% Dieser Teil koennte sonst noch irgendwo untergebracht werden.
%--------
% In general, capture and replay can be divided into two phases: The \emph{capture phase}, where the application is run and the information that is needed for replaying the application is captured and the \emph{replay phase} where the application is replayed based on this information.
% 
% During \emph{capture phase}, the capture and replay implementation needs to record the information that will later be needed to replay the observed part. In general this is at least all information that is passed from the unobserved part to the observed part. The information is captured by some management code that was introduced by the capture and replay framework (\figref{fig:GenericCrStructure_capture}).

% 
% During \emph{replay phase} the management code needs to replace the unobserved part ( 
% \figref{fig:GenericCrStructure_replay}) in order to replay the run of the observed part. Depending on the part of the application that was defined to be unobserved, the management part can act both as a driver (e.g. in the case of mouse events) and stub (e.g. in the case of a network socket). 
% %Schema dazu: Capture phase & Replay phase , observed & unobserved part && log
% \begin{figure}[h]
%   \centering
%   \includegraphics[width=0.4\textwidth]{illustrations/capture_and_replay_generic_structure_replay}
%   \caption{Generic structure of a capture and replay technique - replay phase}
%   \label{fig:GenericCrStructure_replay}
% %\includegraphics{illustrations/capture_and_replay_generic_structure}
% \end{figure}

Different implementations make different assumptions about what is deterministic and what changes its behaviour throughout different runs of the application. Therefore they define different portions of the program as observed and unobserved part. In the following we will categorize the different implementation based on which part is defined to be unobserved.

\section {Capturing User Input}
The best known capture and replay technique is the capturing of user input, which is most often used for GUI testing. Here, keyboard and mouse events are considered to come from the unobserved part. The exact location where these events are captured may vary (inside the application or through the operating system), but the advantages and limitations mostly stay the same. Usually, capture and replay of user input is used for regression testing. In that setup, the tester executes a sequence of actions on the applications GUI while capturing is enabled. These recorded actions can be replayed on the application in order to check if there were regressions. Abbot \cite{abbot} is an example of such a tool, although it offers more features than only capture and replay of user interactions. %TODO: wie funktioniert Abbot?

Compared to our implementation, the capturing of user inputs has many drawbacks due to its simplicity. Because only the user input is considered to belong to the unobserved part, it can replay programs that interact with other parts of the environment (like file system and network) only as long as these interactions are deterministic. One strategy is to restore the environment before every run of the application, but this is not always easy (e.g. if a random generator is used). If the assumption, that the observed space behaves deterministically does not hold, this results in an incorrect replay of the application. Because our implementation makes it possible to individually define the applications observed and unobserved part, it can also replay programs that depend on other parts of the non-deterministic environment than only the user input.

\section {Capturing Interactions with Libraries}
JRapture \cite{jrapture} uses an approach that has the potential of replaying more complex applications. JRapture is a tool for capturing and replaying Java applications in the field. In addition to capturing and replaying, it also offers a profiling interface that permits the program to be instrumented for profiling in the replay phase. JRapture uses the whole Java API api as the unobserved part of the application. The technique is able to capture interaction between the Java API and the rest of the program by manual instrumentation of the Java API classes. JRapture supports multithreaded applications, but it can not guarantee a deterministic replay of concurrent applications.

Unlike our technique, JRapture relies on a manually modified version of libraries. Programs that need to be captured and replayed need to use a special version of the Java API.  This results in the fact that programs that interact with the environment through other mechanisms than the Java API (for example through the Java Native Interface JNI) can not be supported without additional manual instrumentations.

%genauer veranschaulichen, was complete und incomplete ist.
JRapture captures the complete information that is passed between the observed and unobserved part, in contrast to our approach that only captures necessary information.

\section {Capturing Interactions with the Scheduler}
The techniques presented so far did not consider thread scheduling as another source of non-determinism. DejaVu \cite{dejavu}, a capture and replay tool for Java, considers thread scheduling as its only source of non-determinism, thus thread scheduling belongs to its unobserved part. It is not easy to instrument the scheduler in order to detect thread switches, because threads are often scheduled by the operating system. DejaVu therefore introduces the concept of \emph{logical thread schedule} which is a simplified version of the real thread schedule (the \emph{physical thread schedule}). The \emph{logical thread schedule} contains enough thread schedule information to reproduce the execution behaviour of the program under the assumption that the thread schedule is the only source of non-determinism. By detecting some critical events during capture phase such as accesses to shared variables, and synchronization events, DejaVu is able to deduce the logical thread schedule.

DejaVu concentrates on enabling a determistic replay of multi threaded programs that do not have any other source of non-determinism than the thread scheduler. Our technique does not take the non-determinism of the thread scheduler into account, instead it concentrates on all other interactions with the environment. DejaVu and 
%TODO referenz auf Orso.
our technique are orthogonal and could be combined in order to allow replays of general concurrent programs. 

\section{Selective Capture Replay}
In contrast to other techniques, \emph{selective capture and replay} \cite{orso05may} implemented for Java, offers the user the possibility to make his own definition of observed and unobserved part of the system.
%TODO 2x passiv!
Observed and unobserved part are both defined as a set of classes. The interactions between observed and unobserved part are captured using automated code instrumentation.

%TODO: 2 verschiedene implementierungen (???) Zeller vs. Orso

Our work tries to transfer the concept from the original paper \cite{orso05may} to Eiffel Our implementation instruments code in a different way and adds some Eiffel specific modifications.
%TODO: zwei 

%Überblick über das Paper
%Limitations (polymorphism / dynamic binding) --> oder doch nicht? she. Annahme über observed set.
%Benefit for cdd


\section{Capture and Replay for Eiffel}
\subsection{Relevant Differences between Eiffel and Java}
\subsection{Differences in the Requirements}
\subsection{Idea for Code Instrumentation}

% Unterschiedliche Language features Eiffel/Java
% Unterschiede der Anforderungen
% -keine dynamische Kompilierung
%  capture/replay ohne neukompilierung (cdd)

\chapter{Implementation}
\section{Architecture}
%unterschiedliche Parts noch genauer beschreiben
%-Log: (mikro-architektur), wieso kein XML?
\section{Code Instrumentation}
\subsection{Procedures}
\subsection{Direct Manipulation from C Level}
\subsection{Field Accesses}
\section{Advantages over Traditional Capture and Replay}
\section{Limitations}
%unter welchen Umstaenden musste manuell instrumentiert werden??
\section{Future Steps}
%Design des Management Codes
% Idee der Instrumentierung & Beispiel
% Vorteile gegenüber der Original Implementierung
% Limitations
% Experimental Results
\section{Building the Example from Source}
In this section the process of building an example with capture replay support under Linux will be described. Setting up the modified version of Eiffel Studio and all necessary tools will take the biggest part of this explanation.\\
%TODO: wieso genau wird eine modifizierte Version von Eiffel Studio benoetigt??
It is assumed that all commands in the following steps are executed in the same session, thus keeping the environment variables.\\

\subsection{Building the Preliminaries}
The first tools we need are the  Eiffel Studio Tools \cite{estudiotools}. These will be used in many setup scripts in the examples or tests from the repository. Install those tools according to the description on the webpage.\\

The delivery of the modified Eiffel Studio was built using revision 69201 of Eiffel Studio. Building a delivery with a later version of Eiffel Studio was not testet, so it might not work. If no binaries of revision 69201 are available, a delivery of that revision from source \texttt{(stimmt das wirklich??? gaebe das kein Henne / Ei Problem?)}.\\
After copying the delivery to ~/estudio/Eiffel60\_gpl\_69201, it can be activated:
\bashlisting
\begin{lstlisting}
   activate_estudio 60_gpl_69201
\end{lstlisting}

Now Gobo \cite{gobo} can be installed can be installed from svn (Revision 6001 was successfully tested).
- Install Gobo  from svn (revision 6001) -->
\begin{lstlisting}
svn co -r6001 https://gobo-eiffel.svn.sourceforge.net/svnroot/gobo-eiffel/gobo/trunk ~/capture_replay/gobo
export GOBO=~/capture\_replay/gobo
export PATH=$GOBO/bin:$PATH
$GOBO/work/bootstrap/bootstrap.sh gcc ise
\end{lstlisting}

As all preliminaries are installed, Erl-G \cite{erlg} can be downloaded and built. Revision 719 of Erl-G was tested together with capture and replay.
\begin{lstlisting}
svn co -r719 https://svn.origo.ethz.ch/autotest/trunk/erl_g ~/capture_replay/erl_g
export ERL_G=~/capture_replay/erl_g
export PATH=$ERL_G/bin:$PATH
cd $ERL_G
#EIFFEL_SRC is needed. avoid conflicts between EIFFEL_SRC and ISE_LIBRARY.
export ISE_LIBRARY=$EIFFEL_SRC
geant install
geant compile
\title{Selective Capture and Replay for Eiffel
\end{lstlisting}

To build a delivery of the modified Eiffel Studio, execute these commands: (this will take a few hours).
\begin{lstlisting}
cd ~/capture_replay/
mkdir es
svn co https://eiffelsoftware.origo.ethz.ch/svn/es/branches/capture_replay es
export EIFFEL_SRC=~/capture_replay/es/Src
cd es
geant -b $EIFFEL_SRC/scripts/build.eant build_es
\end{lstlisting}

Before an example can be built the delivery that was just created needs to used be set as default instance of Eiffel Studio
\begin{lstlisting}
cd ~/estudio
ln -s ~/capture_replay/es/EiffelXX EiffelCR
activate_estudio CR
\end{lstlisting}

In order to make the created Eiffel Studio use a modified version of the runtime, it is necessary to recompile the runtime with modified CFLAGS. The new version of the runtime then needs to be installed in the delivery.

It is not possible to directly build Eiffel Studio with the modified version of the runtime, because the changes in the runtime are not compatible to Eiffels store mechanism (\texttt{TODO: das auch noch unter irgendwelchem future work anmerken...}). This would render Eiffel Studio unusable because it relies on this mechanism during the build process.
\begin{lstlisting}
export CFLAGS='-DCAPTURE_REPLAY' 
cd $EIFFEL_SRC
#build the runtime from scratch (clobber the old one)
geant -b scripts/build.eant compile_runtime
cd $ISE_EIFFEL/studio/spec/linux-x86/lib
rm *
cp $EIFFEL_SRC/C/run-time/lib* .
\end{lstlisting}


\subsection{Building an Example}
Now, all necessary tools should be installed and the corresponding environment variables set. And we can start to build an example.\\
First we need to add reflection support to the example project. Erl G will generate reflection classes for us. If the environment variables are correctly set, the geant script should invoke Erl G automatically. \\
At the moment Erl-G does not support overrides because this feature is missing in the Gobo parser. Therefore it's necessary to override the necessary classes manually. There are two geant tasks that take care of this:

\begin{itemize}
\item \identifier{patch\_elks} makes the manual override by copying the modified elks classes from \$EIFFEL\_SRC/library/base/capture\_replay/elks\_overrides to \$ISE\_LIBRARY/library/base/elks \\
\item \identifier{unpatch\_elks} restores the original state by copying the original elks classes from \$EIFFEL\_SRC/library/base/elks to \$ISE\_LIBRARY/library/base/elks .
\end{itemize}

\begin{lstlisting}
cd ~/capture_replay/es/examples/capture_replay/iteration1
geant install
\end{lstlisting}

The example can now be opened with the modified version of Eiffel Studio. Make sure that the CFLAGS are still set to '-DCAPTURE\_REPLAY'. Otherwise it will not be possible to build the example.

\section{Experimental Results}
\subsection{Example Application}
\subsection{Results}
\subsection{Conclusion}
%Eiffel Vision Beispiel, mit unterschiedlichen observed sets

\section{Future Work}
%Fehlende C/R und Language Features beschreiben, sowie einen möglichen Lösungsansatz.
\section{Appendix}

\subsection{Log File Grammar}

log ::= event* \\
event ::= call $\mid$ return $\mid$ outread\\
call ::= calltype entity methodname entity* \texttt{“\%N”} \\
return ::= returntype [entity]  \texttt{“\%N”}\\
outread ::= \texttt{OUTREAD} entity attribute\_name entity
calltype ::= \texttt{INCALL} $\mid$ \texttt{OUTCALL} \\
returntype ::= \texttt{INCALLRET} $\mid$ \texttt{OUTCALLRET} \\
methodname ::= identifier \\
attribute\_name ::= identifier \\
entity ::= (object $\mid$ value) \\
size ::= integer \\
object ::= \texttt{[NON$\_$BASIC} typename object\_id \texttt{]} \\
value ::= \texttt{[BASIC} typename \texttt{"} string \texttt{"]} \\
typename ::= identifier \\
object\_id ::= integer\\
identifier ::= [A-Za-z]character*\\
string ::= character*\\

\subsection{capture-phase performance measurements}
\begin{itemize}
	\item normal application: 2.5s
	\item Captured application: 30s
	\begin{itemize}
		\item RECORDER: 1.7s\\
		\item SERIALIZER: 26s
		\begin{itemize}
			\item write$\_$statements: 16s
			\begin{itemize}
				\item file.put$\_$string: 4s
				\item object$\_$id: 7s
				\item other: 5s
			\end{itemize}
			\item is$\_$basic$\_$type: 3s
			\item other: 7s
		\end{itemize}
	\end{itemize}
\end{itemize}

\bibliographystyle{plain} 
\bibliography{thesis}

\end{document}


